%% Page 6: the rest of the ablation (Figure 5) and robustness to the exemplars (Figure 6).
\crumbs{§3.3 Ablation \,›\, Figure~5 \,›\, §3.4 Robustness \,›\, Figure~6}
\insection{ablation,rest}

%% ---- marks on the original page
\bracket{1}{26.7}{108.1}
\badge{2}{32.5}{28.0}
\badge{3}{32.5}{72.3}
\highlight[hlpink]{113.0}{49.4}{125.4}{52.8}
\highlight[hlpink]{38.1}{53.2}{82.4}{56.7}
\highlight[hlpink]{38.1}{97.4}{95.4}{101.0}
\badge{4}{32.5}{126.7}
\bracket{5}{119.6}{218.5}
\bracket{6}{221.1}{232.4}

\block[y=27]{%
\heading{Figure 5: which ingredient matters?}
An \textbf{ablation} swaps out one ingredient at a time to see which one causes the effect.
\num{1}\ GSM8K solve rate for LaMDA 137B and PaLM 540B, five prompt variants: only the last bar
(chain of thought) stands out.}

\block{%
\num{2}\ \textbf{Variable compute only}: the model writes dots (\ldots) as long as the equation, then answers.
It \hl[hlpink]{performs about the same as the baseline}: more tokens alone are not the reason.}

\block[y=71]{%
\num{3}\ \textbf{Chain of thought after the answer}: also about the same as the baseline. The reasoning has to come
\textbf{first} and lead to the answer; it is not just ``activating knowledge''.
\begin{keyidea}
Not the equation, not extra tokens, not reasoning after the fact: the natural-language steps \textbf{before} the
answer are what help.
\end{keyidea}}

\block[y=118]{%
\num{4}\ \textbf{§3.4 Robustness}: prompts are known to be sensitive to the exact examples. So three people
(annotators A, B, C) each wrote their own chains of thought for the same exemplars; B and C did not see A's.}

\block{%
\num{5}\ \textbf{Figure 6} (LaMDA 137B, GSM8K and MAWPS): every chain-of-thought variant beats standard prompting by
a large margin. There is variance between writers, as expected for any prompt, but the gain is never lost.
\begin{example}
Annotator A wrote several short sentences; a deliberately concise style (footnote 1) still works.
\end{example}}

\block[y=221]{%
\num{6}\ Also tested with exemplars taken from GSM8K's own training set: the gains hold, so they are not an
artefact of one writer.}

\block{%
\begin{summary}
Pages 1--6: write the reasoning before the answer, in words, in a few examples, and a large enough model will
reason too.
\end{summary}
\grey{Wei, J. et al. (2022). Chain-of-Thought Prompting Elicits Reasoning in Large Language Models. arXiv:2201.11903,
CC BY 4.0. Original pages reproduced unmodified; marks and notes added.}}
